\documentclass[12pt,a4paper]{article}

\usepackage[utf8]{inputenc}
\usepackage[T1]{fontenc}
\usepackage[polish]{babel}
\usepackage{lmodern}
\usepackage[a4paper,margin=2.5cm]{geometry}
\usepackage{amsmath}
\usepackage{graphicx}
\usepackage{enumitem}
\usepackage{tikz}
\usepackage[hidelinks]{hyperref}

% Grafiki do sprawozdania wrzucać do katalogu zrzuty/
\graphicspath{{zrzuty/}}

\setlist[itemize]{topsep=0.35em,itemsep=0.2em,leftmargin=1.6em}
\setlist[enumerate]{topsep=0.35em,itemsep=0.2em,leftmargin=1.6em}

\newcommand{\kod}[1]{\texttt{#1}}

\begin{document}

% Strona tytułowa jak w strona_tyt_zssk.pdf (Times).
% Bold: nagłówek i temat. Bez bolda: nazwa przedmiotu (24 pt) i etykiety.
\thispagestyle{empty}
\mbox{}%
\begin{tikzpicture}[remember picture, overlay]
	\draw[line width=1.15pt]
		([xshift=97.8pt, yshift=-124.8pt]current page.north west)
		-- ([xshift=528.8pt, yshift=-124.8pt]current page.north west);
	\draw[line width=1.15pt]
		([xshift=97.8pt, yshift=-721.2pt]current page.north west)
		-- ([xshift=528.8pt, yshift=-721.2pt]current page.north west);

	\node[anchor=base, inner sep=0]
		at ([xshift=313.3pt, yshift=-89.5pt]current page.north west)
		{\fontfamily{qtm}\fontsize{20.04}{23.1}\selectfont\bfseries\scshape Politechnika Wrocławska};
	\node[anchor=base, inner sep=0]
		at ([xshift=313.3pt, yshift=-112.6pt]current page.north west)
		{\fontfamily{qtm}\fontsize{20.04}{23.1}\selectfont\bfseries\scshape Wydział Informatyki i Telekomunikacji};

	\node[anchor=base, inner sep=0]
		at ([xshift=313.3pt, yshift=-272.7pt]current page.north west)
		{\fontfamily{qtm}\fontsize{24}{27.6}\selectfont\mdseries\upshape Zarządzanie w systemach i sieciach};
	\node[anchor=base, inner sep=0]
		at ([xshift=313.3pt, yshift=-300.3pt]current page.north west)
		{\fontfamily{qtm}\fontsize{24}{27.6}\selectfont\mdseries\upshape komputerowych - projekt};

	\node[anchor=base, inner sep=0]
		at ([xshift=313.3pt, yshift=-386.4pt]current page.north west)
		{\fontfamily{qtm}\fontsize{20.04}{23.1}\selectfont\bfseries\upshape Mechanizmy zarządzania wątkami programowymi};
	\node[anchor=base, inner sep=0]
		at ([xshift=313.3pt, yshift=-409.5pt]current page.north west)
		{\fontfamily{qtm}\fontsize{20.04}{23.1}\selectfont\bfseries\upshape w zadaniu wykrywania krawędzi na obrazach};
	\node[anchor=base, inner sep=0]
		at ([xshift=313.3pt, yshift=-432.6pt]current page.north west)
		{\fontfamily{qtm}\fontsize{20.04}{23.1}\selectfont\bfseries\upshape cyfrowych z wykorzystaniem operatorów};
	\node[anchor=base, inner sep=0]
		at ([xshift=313.3pt, yshift=-455.7pt]current page.north west)
		{\fontfamily{qtm}\fontsize{20.04}{23.1}\selectfont\bfseries\upshape gradientowych };

	\node[anchor=base west, inner sep=0]
		at ([xshift=99.3pt, yshift=-541pt]current page.north west)
		{\fontfamily{qtm}\fontsize{15.96}{18}\selectfont\mdseries\scshape Autorzy:};
	\node[anchor=base west, inner sep=0]
		at ([xshift=99.3pt, yshift=-563.6pt]current page.north west)
		{\fontfamily{qtm}\fontsize{14.04}{16}\selectfont Karol Gąsior};
	\node[anchor=base west, inner sep=0]
		at ([xshift=99.3pt, yshift=-577.1pt]current page.north west)
		{\fontfamily{qtm}\fontsize{11.04}{13}\selectfont Indeks: 280913};
	\node[anchor=base west, inner sep=0]
		at ([xshift=99.3pt, yshift=-599pt]current page.north west)
		{\fontfamily{qtm}\fontsize{14.04}{16}\selectfont Mateusz Kunicki};
	\node[anchor=base west, inner sep=0]
		at ([xshift=99.3pt, yshift=-612.5pt]current page.north west)
		{\fontfamily{qtm}\fontsize{11.04}{13}\selectfont Indeks: 272468};

	\node[anchor=base west, inner sep=0]
		at ([xshift=287.1pt, yshift=-541pt]current page.north west)
		{\fontfamily{qtm}\fontsize{15.96}{18}\selectfont\mdseries\scshape Prowadzący zajęcia:};
	\node[anchor=base west, inner sep=0]
		at ([xshift=287.1pt, yshift=-563.6pt]current page.north west)
		{\fontfamily{qtm}\fontsize{14.04}{16}\selectfont Dr inż.\ Robert Wójcik, W4N};

	\node[anchor=base west, inner sep=0]
		at ([xshift=287.1pt, yshift=-640.4pt]current page.north west)
		{\fontfamily{qtm}\fontsize{15.96}{18}\selectfont\mdseries\scshape Ocena pracy:};

	\node[anchor=base, inner sep=0]
		at ([xshift=313.3pt, yshift=-742.4pt]current page.north west)
		{\fontfamily{qtm}\fontsize{14.04}{16}\selectfont Wrocław, 2026};
\end{tikzpicture}
\clearpage

%----------------------------------------------------------------------


\section{Temat i cel projektu}

Temat: Mechanizmy zarządzania wątkami programowymi w zadaniu wykrywania krawędzi na obrazach cyfrowych z wykorzystaniem operatorów gradientowych (Sobela i pokrewnych).

Cel projektu: projekt oraz implementacja programu wykrywającego krawędzie na obrazie wczytywanym z dysku, w którym obliczenia rozdzielane są pomiędzy dowolną, zadaną przez użytkownika liczbę wątków. Za podział pracy odpowiadać będzie wbudowany mechanizm szeregujący (\emph{scheduler}) oferujący kilka strategii przydziału fragmentów obrazu. Dodatkowym celem jest zbadanie, jak liczba wątków, strategia szeregowania oraz rozmiar przydzielanego fragmentu wpływają na czas przetwarzania (zegarowy i procesora) oraz na zużycie pamięci operacyjnej.

\section{Sformułowanie problemu}

Problem wykrywania krawędzi będzie rozwiązywany przy uwzględnieniu następujących założeń:
\begin{itemize}
	\item obraz wejściowy wczytywany jest z dysku i od razu konwertowany do skali szarości (8~bitów na piksel);
	\item cały obraz mieści się w pamięci operacyjnej i jest przetwarzany w niej w całości;
	\item użytkownik wybiera jedną z dostępnych masek: Sobel \(3 \times 3\), Sobel \(5 \times 5\), Prewitt \(3 \times 3\), Scharr \(3 \times 3\) oraz krzyż Robertsa \cite{gonzalez,tadeusiewicz};
	\item piksele leżące poza brzegiem obrazu zastępowane są najbliższym pikselem brzegowym (powielenie krawędzi);
	\item wynikiem jest moduł gradientu znormalizowany do zakresu 0--255 i zapisany do pliku \kod{output.jpg};
	\item najmniejszą jednostką podziału pracy pomiędzy wątki jest wiersz obrazu;
	\item liczba wątków, strategia szeregowania i rozmiar fragmentu są parametrami programu;
	\item program działa wyłącznie w trybie tekstowym (terminal), bez interfejsu graficznego.
\end{itemize}

Problem obliczeniowy \cite{gonzalez,tadeusiewicz}: dla obrazu \(I\) o wymiarach \(W \times H\) oraz pary masek \(K_x\), \(K_y\) o rozmiarze \(k \times k\) wyznaczyć dla każdego piksela wartość
\begin{equation}
	G(x, y) = \sqrt{G_x(x, y)^{2} + G_y(x, y)^{2}},
	\quad G_x = K_x * I,
	\quad G_y = K_y * I,
\end{equation}
gdzie \(*\) oznacza dwuwymiarowy splot dyskretny.

Problem badawczy (optymalizacyjny) \cite{williams,openmp}: dobrać liczbę wątków \(N\), strategię szeregowania oraz rozmiar fragmentu tak, aby zapewnić optymalne wartości następujących kryteriów:
\begin{itemize}
	\item czas zegarowy przetwarzania obrazu -- minimum;
	\item efektywność zrównoleglenia \(E(N) = S(N)/N\), gdzie \(S(N) = T(1)/T(N)\) -- maksimum;
	\item nierównomierność obciążenia wątków (różnica czasów zakończenia pracy) -- minimum;
	\item dodatkowe zużycie pamięci RAM związane z wielowątkowością -- minimum.
\end{itemize}

\section{Analiza złożoności obliczeniowej problemu}

Złożoność obliczeniowa splotu wynosi \(O(W \cdot H \cdot k^{2})\), a więc dla ustalonej maski zależy liniowo od liczby pikseli obrazu. Złożoność pamięciowa wynosi \(O(W \cdot H)\) -- program przechowuje obraz wejściowy (1~bajt na piksel), mapę modułu gradientu (4~bajty na piksel) oraz obraz wyjściowy (1~bajt na piksel), czyli około 6~bajtów na piksel.

Wartość każdego piksela wyniku zależy wyłącznie od jego otoczenia w obrazie wejściowym, dlatego zadanie należy do klasy problemów łatwo zrównoleglanych (\emph{embarrassingly parallel}) i teoretycznie umożliwia przyspieszenie liniowe względem liczby rdzeni. W praktyce przyspieszenie ograniczają: część sekwencyjna programu (dekodowanie i kodowanie JPG, redukcja maksimum), narzut synchronizacji oraz przepustowość pamięci. Górne ograniczenie przyspieszenia opisuje prawo Amdahla \cite{amdahl}:
\begin{equation}
	S(N) \le \frac{1}{(1 - p) + \dfrac{p}{N}},
\end{equation}
gdzie \(p\) oznacza udział części programu, którą można wykonywać równolegle.

\section{Metoda i algorytmy rozwiązywania problemu}

Zastosowana zostanie dekompozycja danych: obraz dzielony jest na fragmenty złożone z kolejnych wierszy, a każdy wątek przetwarza otrzymane fragmenty niezależnie, zapisując wyniki do rozłącznych obszarów pamięci, co eliminuje potrzebę stosowania blokad. Przetwarzanie przebiega w dwóch etapach rozdzielonych barierą:
\begin{enumerate}
	\item splot z maską oraz wyznaczenie lokalnego maksimum gradientu w każdym wątku;
	\item redukcja lokalnych maksimów do maksimum globalnego, a następnie równoległa normalizacja wyniku do zakresu 0--255.
\end{enumerate}

Wątki tworzone są jednokrotnie i utrzymywane w puli (\emph{thread pool}), dzięki czemu kolejne etapy nie ponoszą kosztu ich ponownego tworzenia. Przydziałem fragmentów obrazu zarządza scheduler oferujący strategie wzorowane na klauzuli \emph{schedule} standardu OpenMP \cite{openmp}:
\begin{itemize}
	\item \emph{static} -- obraz dzielony na \(N\) równych bloków, po jednym na wątek;
	\item \emph{cyclic} -- fragmenty o stałej wielkości przydzielane cyklicznie (round-robin);
	\item \emph{dynamic} -- wolny wątek pobiera kolejny fragment ze wspólnej kolejki (licznik atomowy);
	\item \emph{guided} -- jak \emph{dynamic}, lecz rozmiar fragmentów maleje wraz z ilością pozostałej pracy.
\end{itemize}

\section{Metoda, technologie i narzędzia implementacji}

Program zostanie napisany w języku C++ (standard C++17) z wykorzystaniem biblioteki STL, w szczególności mechanizmów wielowątkowości: \kod{std::thread}, \kod{std::mutex}, \kod{std::condition\_variable} oraz \kod{std::atomic} \cite{williams,cppreference}. Odczyt i zapis plików JPG zrealizowany zostanie przy użyciu jednoplikowych bibliotek \kod{stb\_image} i \kod{stb\_image\_write} \cite{stb}. Projekt będzie budowany za pomocą CMake i realizowany w środowisku Qt Creator (bez wykorzystania bibliotek Qt). Logika przetwarzania zostanie zamknięta w klasie \kod{EdgeDetector}, której parametrami są liczba wątków, strategia szeregowania oraz rozmiar fragmentu.

\section{Sposób testowania i oceny jakości rozwiązań}

Poprawność działania zostanie zweryfikowana poprzez porównanie wyników wersji wielowątkowej z wynikiem wersji jednowątkowej -- obrazy wyjściowe muszą być identyczne niezależnie od liczby wątków i strategii szeregowania. Testy przeprowadzone zostaną na obrazach syntetycznych o znanym położeniu krawędzi oraz na zdjęciach rzeczywistych.

Ocena wydajności obejmie pomiar dla każdego etapu (wczytanie, splot, normalizacja, zapis): czasu zegarowego -- przy użyciu biblioteki \kod{chrono} STL \cite{cppreference}, czasu procesora całego procesu oraz poszczególnych wątków, a także bieżącego i szczytowego zużycia pamięci RAM (funkcje systemowe Windows i Linux). Wbudowany tryb porównawczy wykona przetwarzanie dla \(1, 2, 4, \ldots, N\) wątków z kilkukrotnym powtórzeniem pomiaru i wyznaczy przyspieszenie oraz efektywność. Badania zostaną przeprowadzone dla obrazów o różnej rozdzielczości (od HD do 8K), różnych masek, strategii szeregowania i rozmiarów fragmentu.

W aplikacji nie przewiduje się interfejsu graficznego. Wyniki pomiarów będą wypisywane w terminalu oraz zapisywane do plików tekstowych, a następnie analizowane. W celu wizualnej prezentacji wyników badań przedstawione zostaną odpowiednie wykresy.

% Przykład wstawienia grafiki z katalogu zrzuty/ (odkomentować po dodaniu pliku):
% \begin{figure}[ht]
%   \centering
%   \includegraphics[width=0.85\textwidth]{nazwa.png}
%   \caption{Opis rysunku.}
% \end{figure}

\begin{thebibliography}{9}
	\bibitem{gonzalez}
	Gonzalez R.~C., Woods R.~E., \emph{Digital Image Processing}, 4th ed., Pearson, 2018.

	\bibitem{tadeusiewicz}
	Tadeusiewicz R., Korohoda P., \emph{Komputerowa analiza i przetwarzanie obrazów}, Wydawnictwo Fundacji Postępu Telekomunikacji, Kraków, 1997.

	\bibitem{williams}
	Williams A., \emph{C++ Concurrency in Action}, 2nd ed., Manning, 2019.

	\bibitem{openmp}
	OpenMP Architecture Review Board, \emph{OpenMP Application Programming Interface}, Version 5.2, 2021.

	\bibitem{amdahl}
	Amdahl G.~M., \emph{Validity of the single processor approach to achieving large scale computing capabilities}, AFIPS Conference Proceedings, vol.~30, 1967.

	\bibitem{cppreference}
	\url{https://en.cppreference.com/w/cpp/thread}, 19.10.2026.

	\bibitem{stb}
	\url{https://github.com/nothings/stb}, 19.10.2026.
\end{thebibliography}

\end{document}
